% !TEX TS-program = xelatex
\documentclass[12pt,a4paper]{article}

\usepackage{graphicx}
\usepackage{float}
\usepackage{geometry}
\geometry{left=20mm,right=20mm,top=22mm,bottom=22mm}
\usepackage{booktabs}
\usepackage{longtable}
\usepackage{array}
\usepackage{caption}
\usepackage{subcaption}
\usepackage{listings}
\lstset{basicstyle=\ttfamily\small,breaklines=true}
\usepackage{placeins}
\usepackage{amsmath}
\usepackage{amssymb}
\usepackage{mathtools}
\usepackage{hyperref}
\hypersetup{colorlinks=true, linkcolor=blue, urlcolor=blue}

\usepackage{xepersian}
\settextfont[Path=fonts/, Extension=.ttf]{BNazanin}
\setlatintextfont{Times New Roman}
\newcommand{\EN}[1]{\lr{\latinfont #1}}
\newcommand{\ENp}[1]{(\EN{#1})}
\newcolumntype{P}[1]{>{\raggedright\arraybackslash}p{#1}}

\setlength{\parskip}{6pt}
\setlength{\parindent}{0pt}
\renewcommand{\labelitemi}{-}
\setlength{\tabcolsep}{3pt}
\setlength{\LTleft}{0pt}
\setlength{\LTright}{0pt}

\begin{document}
\begin{latin}
\section{Database and Data Construction Workflow}
This chapter documents the full data pipeline used to build the final modeling table, from Google Earth Engine extraction to the final labeled and feature-engineered dataset.

\subsection{Pipeline Summary}
The final implemented chain is:
\begin{enumerate}
\item \texttt{drive/DataLine\_(1).ipynb}: final Earth Engine extraction at \texttt{3000 m} for a \texttt{20 x 20} grid, split into fast-changing meteorology (\texttt{MET}) and slower bio-climate (\texttt{BIO}) streams over 2000--2020.
\item \texttt{drive/DataLine\_(1).ipynb}: monthly CSV export for each stream, then inner merge on \texttt{(node,date)} to build \texttt{GRID3km\_DAILY\_20x20\_merged.csv}.
\item \texttt{drive/DataLine\_(1).ipynb}: final rescaling and cleanup (NDVI/Albedo/AAT scaling, optional LAI drop) to save \texttt{GRID3km\_DAILY\_20x20\_modelready.csv}.
\item \texttt{drive/Labeling\_grid3km.ipynb}: VISI-based label construction and merge to produce \texttt{GRID3km\_DAILY\_20x20\_modelready\_labeled.csv}.
\item \texttt{Build Data2/new/run.py}: feature engineering (anomalies, rolling windows, wind-direction, erosion, neighborhood contrasts) from labeled model-ready data.
\item \texttt{convert.py} (or equivalent parquet-to-CSV step): conversion to the final file \texttt{GRID3km\_DAILY\_20x20\_modelready\_labeled\_featured.csv}.
\end{enumerate}

\subsection{Earth Engine Extraction Setup (Final Run)}
The extraction setup in \texttt{drive/DataLine\_(1).ipynb} is:
\begin{itemize}
\item Project: \texttt{starter-458908}.
\item Region: rectangle \texttt{[45.389892, 31.230722, 45.928881, 31.690477]} (W, S, E, N).
\item CRS: \texttt{EPSG:4326}.
\item Sampling scale: \texttt{SCALE = 3000 m}.
\item Grid size: \texttt{20 x 20} nodes (named \texttt{i\_j}).
\item Time range: \texttt{2000-01-01} to \texttt{2020-12-31} (20 years).
\item MET stream: ERA5 and MERRA are aggregated to daily means.
\item BIO stream: MODIS/TerraClimate/SPEI features are sampled with hold-last windows.
\item Missing strategy: fill value \texttt{-999} is used before later cleaning.
\item Sampling method: \texttt{sampleRegions(..., scale=SCALE, projection=EPSG:4326, tileScale=4)}.
\end{itemize}

\noindent\textbf{Resolution note for the final modeling table.}
The final dataset name \texttt{GRID3km\_DAILY\_20x20\_modelready...} refers to the downstream node grid used in modeling and labeling.
With bounds \texttt{W=45.389892, S=31.230722, E=45.928881, N=31.690477} and \texttt{NX=NY=20}, the node spacing is approximately
\texttt{0.02695 deg} (lon) and \texttt{0.02299 deg} (lat), i.e., about \texttt{2.56 km x 2.56 km} near the study latitude (commonly rounded and referred to as ``3 km'').

\subsection{Earth Engine Feature Catalog Used in Final Dataset}

\begin{latin}

{\small
\begin{longtable}{@{}P{2.2cm} P{4.2cm} P{3.0cm} P{6.0cm}@{}}
\toprule
\textbf{Feature} & \textbf{Earth Engine Asset ID} & \textbf{Band / Formula} & \textbf{Description} \\
\midrule
\endfirsthead
\toprule
\textbf{Feature} & \textbf{Earth Engine Asset ID} & \textbf{Band / Formula} & \textbf{Description} \\
\midrule
\endhead
\midrule
\multicolumn{4}{r}{Continued on next page} \\
\endfoot
\bottomrule
\endlastfoot

u10 & \texttt{ECMWF/ERA5\_LAND/ HOURLY} & \texttt{u\_comp...} & Eastward component of wind at 10m height. \\ \addlinespace

v10 & \texttt{ECMWF/ERA5\_LAND/ HOURLY} & \texttt{v\_comp...} & Northward component of wind at 10m height. \\ \addlinespace

wind\_speed & \texttt{ECMWF/ERA5\_LAND/ HOURLY} & $\sqrt{u10^2 + v10^2}$ & Total horizontal wind velocity derived from u and v components. \\ \addlinespace

merra\_dust & \texttt{NASA/GSFC/MERRA/ aer/2DUCMASS} & \texttt{DUCMASS} & Total dust column mass density; indicates atmospheric aerosol loading. \\ \addlinespace

NDVI & \texttt{MODIS/061/MOD13Q1} & \texttt{NDVI} & Normalized Difference Vegetation Index; measures vegetation greenness. \\ \addlinespace

SAVI & \texttt{MODIS/061/MOD09A1} & Formula* & Soil Adjusted Vegetation Index; minimizes soil brightness influences. \\ \addlinespace

LAI & \texttt{MODIS/061/MOD15A2H} & \texttt{Lai\_500m} & Leaf Area Index; quantifies the amount of leaf material in the canopy. \\ \addlinespace

Precip & \texttt{IDAHO\_EPSCOR/ TERRACLIMATE} & \texttt{pr} & Accumulated precipitation; used to monitor moisture availability. \\ \addlinespace

Evapotransp & \texttt{IDAHO\_EPSCOR/ TERRACLIMATE} & \texttt{aet} & Actual Evapotranspiration; water flux from land to atmosphere. \\ \addlinespace

AAT & \texttt{IDAHO\_EPSCOR/ TERRACLIMATE} & \texttt{(tmmn+tmmx)/2} & Average Air Temperature; mean of daily minimum and maximum temps. \\ \addlinespace

SPEI3 & \texttt{CSIC/SPEI/2\_10} & \texttt{SPEI\_03\_month} & Standardized Precipitation-Evapotranspiration Index (3-month scale). \\ \addlinespace

Albedo & \texttt{MODIS/061/MCD43A3} & \texttt{Albedo\_WSA...} & Shortwave bi-hemispherical reflectance; measures surface reflectivity. \\

\end{longtable}
}

\noindent\small{*SAVI Formula: $1.5 \cdot \frac{NIR-RED}{NIR+RED+0.5}$}

\end{latin}
\noindent\textbf{Important:} variables from earlier experiments such as \texttt{Nightlights}, \texttt{PBLH}, \texttt{VOD}, \texttt{BSCI}, \texttt{ATI}, and \texttt{Delta\_TS} are not part of the final \texttt{GRID3km\_DAILY\_20x20\_modelready\_labeled\_featured.csv} schema.

\subsection{Fast/Slow Split and Merge}
The 20-year extraction is executed as two streams:
\begin{itemize}
\item MET stream (\texttt{u10, v10, wind\_speed, merra\_dust}), exported monthly as 

\texttt{GRID3km\_DAILY\_MET\_YYYYMM\_20x20.csv}.
\item BIO stream (\texttt{NDVI, SAVI, LAI, Precip, Evapotransp, AAT, SPEI3, Albedo}), exported monthly as \texttt{GRID3km\_DAILY\_BIO\_YYYYMM\_20x20.csv}.
\end{itemize}
Then, in \texttt{drive/DataLine\_(1).ipynb}, both streams are read and merged by:
\begin{enumerate}
\item reading all monthly files from Drive,
\item normalizing \texttt{node} and \texttt{date} types,
\item dropping accidental stream-level duplicates on \texttt{(node,date)},
\item inner merge on \texttt{(node,date)} with \texttt{validate='one\_to\_one'},
\item replacing \texttt{-999} sentinels with \texttt{NaN},
\item saving \texttt{GRID3km\_DAILY\_20x20\_merged.csv}.
\end{enumerate}

\subsection{Model-Ready Construction}
The merged table is converted to model-ready format in \texttt{drive/DataLine\_(1).ipynb}:
\begin{itemize}
\item \texttt{Albedo} cleaned to valid range and scaled to unitless (\texttt{*0.001}),
\item \texttt{NDVI} scaled (\texttt{*1e-4}) and clipped to \texttt{[-1,1]},
\item \texttt{AAT} scaled from \texttt{0.1$^\circ$C} units (\texttt{*0.1}),
\item optional monthly-to-daily conversion for \texttt{Precip}/\texttt{Evapotransp} is documented,
\item \texttt{LAI} is dropped in the final model-ready table.
\end{itemize}
Output: \texttt{GRID3km\_DAILY\_20x20\_modelready.csv}.

\subsection{Label Construction and Merge}
After producing \texttt{GRID3km\_DAILY\_20x20\_modelready.csv}, labels are created in 

\texttt{drive/Labeling\_grid3km.ipynb} using VISI events:
\begin{enumerate}
\item parse VISI geometry (\texttt{POINT(lon lat)}),
\item keep points inside the same 20x20 domain,
\item build node centroids and index with a haversine \texttt{BallTree},
\item assign VISI events to neighboring nodes using radius \texttt{R=5 km},
\item create unique \texttt{(node,date)} positives and left-merge to the grid table,
\item set missing labels to zero.
\end{enumerate}
This produces \texttt{GRID3km\_DAILY\_20x20\_modelready\_labeled.csv}.

\subsection{Feature Engineering to Final Dataset}
In \texttt{Build Data2/new/run.py}, the labeled table is transformed into the final featured dataset:
\begin{enumerate}
\item sort by \texttt{node,date} and compute day-of-year climatology for base drivers,
\item compute standardized anomalies (\texttt{wind\_speed\_anom}, \texttt{NDVI\_anom}, \texttt{SPEI3\_anom}, \texttt{AAT\_anom}, \texttt{Albedo\_anom}, \texttt{merra\_dust\_anom}),
\item add rolling temporal context (\texttt{wind\_speed\_max\_3d}, \texttt{wind\_speed\_mean\_7d}, \texttt{Precip\_sum\_30d}, \texttt{Evapo\_sum\_30d}, \texttt{water\_deficit\_30d}, \texttt{NDVI\_mean\_30d}, \texttt{NDVI\_trend\_30d}, \texttt{days\_since\_rain}),
\item derive wind-direction cyclic terms (\texttt{wind\_dir\_sin}, \texttt{wind\_dir\_cos}),
\item derive erosion indices (\texttt{erosion\_index\_1}, \texttt{erosion\_index\_2}),
\item compute 4-neighbor spatial means and contrasts 

(\texttt{NDVI\_nb\_mean}, \texttt{Albedo\_nb\_mean}, \texttt{wind\_speed\_nb\_mean}, \texttt{NDVI\_contrast}, 

\texttt{Albedo\_contrast}, \texttt{wind\_speed\_contrast}),
\item save featured parquet, then convert to final CSV.
\end{enumerate}

\subsection{Final Column Schema (Exact)}
\begin{lstlisting}[basicstyle=\ttfamily\scriptsize,breaklines=true]
['node', 'date', 'label', 'u10', 'v10', 'wind_speed', 'merra_dust',
 'NDVI', 'SAVI', 'Precip', 'Evapotransp', 'AAT', 'SPEI3', 'Albedo',
 'wind_speed_anom', 'NDVI_anom', 'SPEI3_anom', 'AAT_anom', 'Albedo_anom',
 'merra_dust_anom', 'wind_speed_max_3d', 'wind_speed_mean_7d',
 'Precip_sum_30d', 'Evapo_sum_30d', 'water_deficit_30d', 'NDVI_mean_30d',
 'NDVI_trend_30d', 'days_since_rain', 'wind_dir_sin', 'wind_dir_cos',
 'erosion_index_1', 'erosion_index_2', 'NDVI_nb_mean', 'Albedo_nb_mean',
 'wind_speed_nb_mean', 'NDVI_contrast', 'Albedo_contrast',
 'wind_speed_contrast']
\end{lstlisting}

\subsection{Final Dataset Meaning}
The final file \texttt{GRID3km\_DAILY\_20x20\_modelready\_labeled\_featured.csv} is a node-date supervised table where:
\begin{itemize}
\item each row is one grid node at one date,
\item \texttt{label=1} indicates a mapped VISI event day for that node,
\item predictors combine base Earth Engine drivers with temporal anomaly, rolling-window, and neighborhood-contrast features.
\end{itemize}
This structure is designed for rare-event dust prediction with both temporal and spatial context.

\section{GTG Method and Comparative Modeling Results}

\subsection{GTG Idea: Long-Horizon Control with Grouped Temporal Gates}
The \texttt{GTG-TGCN} model addresses a key problem in dust early-warning: features evolve at very different temporal scales.
Instead of forcing a single lookback for all variables, GTG splits features into temporal groups (fast/medium/slow) and learns separate temporal importance for each group.

In our implementation (\texttt{gtg.py}), grouped temporal memory is defined as:
\begin{itemize}
\item fast group: \texttt{14} daily lags (\texttt{stride=1}),
\item medium group: \texttt{112} lag slots with \texttt{stride=7} (weekly spacing),
\item slow group: \texttt{36} lag slots with \texttt{stride=30} (monthly spacing).
\end{itemize}

This design gives a large effective historical horizon while keeping parameter count and memory manageable.
Operationally, GTG computes a group-wise scalar representation from lagged node histories, then the graph backbone (\texttt{GCN}-based in this variant) performs spatial-temporal propagation for final anomaly probability estimation.

\subsection{Explainability Advantages of GTG}
GTG is directly interpretable at feature-lag level.
For each alerted node-date, the explanation module returns top contributors with:
\begin{itemize}
\item feature name,
\item temporal group (fast/med/slow),
\item lag in days and corresponding lag date,
\item learned attention weight,
\item contribution magnitude.
\end{itemize}

This supports answers to operational questions such as ``why this alarm now?'' and ``which historical signals drove the decision?'', while still preserving graph-based spatial context.

\subsection{Comparison with Baseline Models}
To compare fairly, baseline models were trained on the same final dataset and evaluated with the same metrics and the same 5-day early-warning definition used for GTG.
The table below reports test metrics.

\begin{table}[H]
\centering
\caption{Test-set comparison: best GTG run vs baseline models}
\label{tab:model_compare}
\begin{tabular}{lccccc}
\toprule
Model & ROC-AUC & PR-AUC & F1 & Hit\_Rate\_5d & FA\_Rate \\
\midrule
GTG-TGCN (best, seed 77) & 0.7273 & 0.0134 & 0.0263 & 0.924 & 0.13046 \\
gcn & 0.8667 & 0.0334 & 0.0681 & 0.929 & 0.22580 \\
gat & 0.8753 & 0.0359 & 0.0676 & 0.915 & 0.18981 \\
gdn & 0.8558 & 0.0333 & 0.0701 & 0.888 & 0.21188 \\
cnn & 0.8006 & 0.0192 & 0.0454 & 0.895 & 0.32944 \\
lstm & 0.8385 & 0.0311 & 0.0790 & 0.913 & 0.30754 \\
\bottomrule
\end{tabular}
\end{table}

\noindent\textbf{Interpretation.}
On threshold-independent classification metrics (ROC-AUC/PR-AUC), several baselines score higher.
However, for early-warning operation, GTG provides a stronger alarm-quality tradeoff in this experiment, achieving high event hit rate with substantially lower false-alarm rate than the listed baselines.
Because deployment cost is highly sensitive to false alarms, this FA-aware behavior is a primary reason GTG remains the preferred model for the next RL enhancement stage.

\section{Reinforcement Learning Enhancement on Top of GTG}

\subsection{RL Formulation}
The RL stage is applied after GTG and uses GTG output probability as a calibrated prior signal.
For each node-date sample, the agent observes:
\begin{itemize}
\item \texttt{state} = \texttt{[p\_GTG, x\_1, ..., x\_F]},
\item \texttt{action} in \{0,1\} where 1 means ``raise alert''.
\end{itemize}

The reward is asymmetric to strongly penalize false alarms and missed events, with additional alert cost:
\begin{itemize}
\item \texttt{TP}=+8,\ \texttt{FN}=-35,\ \texttt{FP}=-70,\ \texttt{TN}=0,
\item \texttt{alert\_cost}=-1.5 (applied on alert actions for negatives).
\end{itemize}
This converts the final warning decision into a cost-aware policy optimization problem instead of fixed-threshold classification.

\subsection{Training and Evaluation Setup}
The policy is trained using \texttt{PPO} in \texttt{train\_gtg\_rl.py} on \texttt{gtg\_s77\_rl\_data.npz}.
The RL region is the last 8 years of the timeline (with \texttt{total\_years=20}, \texttt{gtg\_train\_years=12}), then split by time into \texttt{70\%/15\%/15\%} for RL train/val/test.
Threshold for operational alarm is selected on RL validation data using the same 5-day early-warning criterion.

\subsection{RL Results}
Table~\ref{tab:rl_results} reports two operating points on the same RL test split:
\begin{itemize}
\item validation-selected threshold (\texttt{thr=0.900}),
\item a nearby operating point targeting approximately \texttt{6\%} false-alarm rate (\texttt{thr=0.869}).
\end{itemize}

\begin{table}[H]
\centering
\caption{RL policy performance on GTG-based inputs (test split)}
\label{tab:rl_results}
\begin{tabular}{lccccccccc}
\toprule
Setting & Thr & Acc & Prec & Rec & F1 & Hit\_5d & Miss\_5d & FA\_Rate & Mean Lead (d) \\
\midrule
RL (val-selected) & 0.900 & 0.9484 & 0.0425 & 0.3118 & 0.0748 & 0.842 & 0.158 & 0.04736 & 3.02 \\
RL (\~6\% FA point) & 0.869 & 0.9364 & 0.0416 & 0.3860 & 0.0751 & 0.879 & 0.121 & 0.05993 & 3.21 \\
\bottomrule
\end{tabular}
\end{table}

\noindent\textbf{Interpretation.}
Compared with GTG-only operation (Hit\_5d=\texttt{0.924}, FA\_Rate=\texttt{0.13046}), RL substantially reduces false alarms (down to \texttt{0.04736} at \texttt{thr=0.900}), but with reduced hit rate.
When threshold is relaxed to \texttt{0.869}, hit rate improves to \texttt{0.879} while FA remains near \texttt{6\%}.
Therefore, RL provides a controllable operating frontier between miss rate and false alarms, and threshold selection should follow deployment cost priorities.

\section{Explainability Cases and Parameter-Interaction Discovery}

\subsection{Three Anomalous Case Visualizations}
To move from prediction to ecosystem understanding, we first inspected individual true-positive explanations.
Figures~\ref{fig:anom_case1}--\ref{fig:anom_case3} show three events from \texttt{plots\_all\_anom}.

\begin{figure}[H]
\centering
\includegraphics[width=0.95\textwidth]{../plots_all_anom/nid1767422_node2_11_event2012-06-05.png}
\caption{Case 1: event 2012-06-05 at node 2\_11.}
\label{fig:anom_case1}
\end{figure}

\begin{figure}[H]
\centering
\includegraphics[width=0.95\textwidth]{../plots_all_anom/nid1795982_node2_19_event2012-08-15.png}
\caption{Case 2: event 2012-08-15 at node 2\_19.}
\label{fig:anom_case2}
\end{figure}

\begin{figure}[H]
\centering
\includegraphics[width=0.95\textwidth]{../plots_all_anom/nid1796044_node4_2_event2012-08-16.png}
\caption{Case 3: event 2012-08-16 at node 4\_2.}
\label{fig:anom_case3}
\end{figure}

\noindent\textbf{Observed pattern from these examples.}
The three cases consistently show a short-term wind activation pattern (\texttt{wind\_speed}, \texttt{wind\_speed\_max\_3d}, \texttt{wind\_speed\_mean\_7d}, lags around 0--13 days) co-occurring with slower background terms (for example \texttt{AAT\_anom} around \texttt{273} days and slow hydro-climate terms around yearly-to-multi-year lags).
This supports the GTG design assumption that dust activation is typically a multi-timescale process rather than a single-day trigger.

\subsection{Method for Mining Cross-Parameter Relations}
After expanding explanations to all test positives (\texttt{1441} events, \texttt{11528} explanation rows), we mined interaction patterns using \texttt{mine\_parameter\_interactions.py}.
The procedure is:
\begin{enumerate}
\item keep strongest contribution per \texttt{(event, feature)} to avoid repeated same-feature rows per event,
\item keep top-\texttt{k} features per event (\texttt{k=8}),
\item discretize lag into interpretable bins (\texttt{0--7}, \texttt{7--30}, \texttt{30--90}, \texttt{90--180}, \texttt{180--365}, \texttt{365--730}, \texttt{730--1095}, \texttt{>1095}),
\item generate all pair combinations inside each event,
\item aggregate across events and rank pairs by frequency, contribution strength, and association (\texttt{lift}, composite \texttt{pair\_score}).
\end{enumerate}

This gives statistical association patterns between lagged parameters; it is evidence of co-occurrence structure, not strict causal proof.

\subsection{Top Interaction Relations}
Table~\ref{tab:interaction_relations} reports top stable relations from \texttt{analysis\_out\_rich/interaction\_pairs\_shortlist.csv}.

\begin{table}[H]
\centering
\caption{Top mined feature-lag interaction relations (test-positive explanations)}
\label{tab:interaction_relations}
\small
\begin{tabular}{P{3.9cm}P{3.9cm}ccccc}
\toprule
Relation A & Relation B & Events & Support & Mean Sum\_Abs & Lift & Both Anom. Rate \\
\midrule
\texttt{Precip} (730--1095d) & \texttt{erosion\_index\_1} (0--7d) & 53 & 0.0368 & 2.089 & 8.65 & 0.151 \\
\texttt{AAT} (365--730d) & \texttt{erosion\_index\_1} (0--7d) & 51 & 0.0354 & 2.070 & 5.37 & 0.000 \\
\texttt{SPEI3} (365--730d) & \texttt{wind\_speed\_mean\_7d} (0--7d) & 113 & 0.0784 & 2.979 & 1.59 & 0.920 \\
\texttt{SPEI3} (365--730d) & \texttt{wind\_speed\_mean\_7d} (7--30d) & 66 & 0.0458 & 2.512 & 2.17 & 0.833 \\
\texttt{SPEI3} (365--730d) & \texttt{SPEI3\_anom} (365--730d) & 376 & 0.2609 & 1.904 & 1.14 & 0.191 \\
\texttt{SPEI3} (365--730d) & \texttt{wind\_speed} (0--7d) & 225 & 0.1561 & 1.956 & 1.26 & 0.533 \\
\texttt{SPEI3} (365--730d) & \texttt{wind\_speed\_max\_3d} (0--7d) & 62 & 0.0430 & 2.165 & 1.65 & 0.710 \\
\texttt{SPEI3} (365--730d) & \texttt{u10} (7--30d) & 79 & 0.0548 & 2.042 & 1.64 & 0.101 \\
\bottomrule
\end{tabular}
\end{table}

\noindent\textbf{Interpretation.}
The dominant structure is a repeated coupling of slow

(\texttt{SPEI3}/\texttt{Precip}/\texttt{AAT} in long lags) with near-term wind and erosion indicators.
Operationally, this indicates that the strongest alarms emerge when long-horizon background stress and short-horizon transport/erosion activation align in time.

\subsection{Robustness Validation of Interaction Claims}
To strengthen scientific claims, we ran three additional validations using \texttt{validate\_interaction\_claims.py} on three seeds (\texttt{s11}, \texttt{s42}, \texttt{s77}).

\begin{table}[H]
\centering
\caption{Robustness checks for mined interactions}
\label{tab:interaction_validation_summary}
\begin{tabular}{P{5.1cm}P{8.8cm}}
\toprule
Check & Result \\
\midrule
Seed-to-seed stability (Top-50 pairs) & Mean pairwise Jaccard=\texttt{0.0689}, median=\texttt{0.0753}; pairs present in all 3 seeds: \texttt{1}; pairs present in at least 2 seeds: \texttt{17}. \\
Permutation significance (reference seed \texttt{s77}, top-30 pairs, 500 permutations) & Significant at \texttt{p<0.05}: \texttt{26/30}; significant at \texttt{FDR<0.10}: \texttt{27/30}. \\
External temporal holdout (split date \texttt{2012-06-10}) & Discovery events: \texttt{1022}, holdout events: \texttt{419}; confirmed pairs (minimum 5 holdout events): \texttt{8/30} (\texttt{26.7\%}). \\
\bottomrule
\end{tabular}
\end{table}

\begin{table}[H]
\centering
\caption{Interactions supported by holdout confirmation and permutation significance}
\label{tab:interaction_multicheck}
\small
\begin{tabular}{P{4.3cm}P{4.1cm}ccccc}
\toprule
Relation A & Relation B & Disc. Events & Holdout Events & Holdout Lift \\
\midrule
\texttt{erosion\_index\_2} (7--30d) & \texttt{wind\_speed\_max\_3d} (7--30d) & 38 & 37 & 2.686 \\
\texttt{SPEI3} (365--730d) & \texttt{wind\_speed\_nb\_mean} (0--7d) & 106 & 23 & 2.793 \\
\texttt{AAT\_anom} (180--365d) & \texttt{SPEI3} (365--730d) & 217 & 46 & 1.416 \\
\texttt{SPEI3} (365--730d) & \texttt{wind\_speed} (0--7d) & 187 & 38 & 2.292 \\
\texttt{SPEI3} (365--730d) & \texttt{v10} (0--7d) & 116 & 19 & 1.714 \\
\bottomrule
\end{tabular}
\end{table}

\noindent\textbf{Validation interpretation.}
The interaction structure is statistically non-random (strong permutation significance), and a subset generalizes to later-time holdout data.
However, seed-to-seed top-list overlap is limited, so conclusions should be framed as robust association hypotheses rather than definitive causal pathways.

\section{Final Statistical Conclusion}
\noindent\textbf{End-to-end outcome.}
This study built a full pipeline from Earth Engine-derived node-date data to operational early warning, then to policy optimization and ecosystem-relationship analysis.
The statistical evidence supports three main conclusions:
\begin{enumerate}
\item \textbf{GTG is operationally competitive.}  
Compared to baseline architectures, GTG remains strong under the early-warning objective, with high event detection and controlled false alarms in the selected operating point (\texttt{Hit\_5d=0.924}, \texttt{FA\_Rate=0.13046}).

\item \textbf{RL provides a controllable alarm-quality frontier.}  
Using GTG outputs as state priors, PPO significantly reduced false alarms at stricter thresholds (\texttt{FA\_Rate=0.04736} at \texttt{thr=0.900}) with the expected hit-rate reduction (\texttt{Hit\_5d=0.842}); a relaxed threshold (\texttt{thr=0.869}) improved hit rate (\texttt{0.879}) while keeping FA near \texttt{6\%}.

\item \textbf{Interaction mining reveals structured multi-timescale associations.}  
From \texttt{1441} explained positive events and \texttt{11528} feature-lag explanation rows, mined relations repeatedly linked slow hydro-climate memory terms (\texttt{SPEI3}, \texttt{AAT}, \texttt{Precip}) with short-term wind/erosion activation terms.
Additional validation showed strong non-randomness (permutation significance: \texttt{26/30} at \texttt{p<0.05}, \texttt{27/30} at \texttt{FDR<0.10}) and partial forward-time confirmation (\texttt{8/30} confirmed in holdout), while seed-to-seed overlap remained limited (mean Jaccard@50=\texttt{0.0689}).
\end{enumerate}

\noindent\textbf{Final claim.}
The presented framework is statistically solid for real-time warning and for generating testable ecosystem interaction hypotheses.
The predictive and RL stages are mature enough for deployment-oriented experiments, and the interaction stage is strong for hypothesis generation, with the caveat that causal interpretation requires further domain-controlled validation.
\end{latin}

\end{document}
